\chapter{Literatura recente e lacuna candidata}\label{literatura-recente-que-define-a-oportunidade}

{
\begingroup
\small
\singlespacing
\setlength{\tabcolsep}{4pt}
\begin{longtable}[]{@{}
  >{\raggedright\arraybackslash}p{(\linewidth - 4\tabcolsep) * \real{0.25}}
  >{\raggedright\arraybackslash}p{(\linewidth - 4\tabcolsep) * \real{0.375}}
  >{\raggedright\arraybackslash}p{(\linewidth - 4\tabcolsep) * \real{0.375}}@{}}
\toprule\noalign{}
\begin{minipage}[b]{\linewidth}\raggedright
Trabalho
\end{minipage} & \begin{minipage}[b]{\linewidth}\raggedright
Resultado ou alcance relevante
\end{minipage} & \begin{minipage}[b]{\linewidth}\raggedright
Consequência para a proposta
\end{minipage} \\ \addlinespace[4pt]
\midrule\noalign{}
\endhead
\bottomrule\noalign{}
\endlastfoot
\autocite{wu2024} & Analisa massa do gráviton com NANOGrav 15 anos.
Distingue a comparação entre modelos da informação associada à
frequência mínima. & Separar informação dos dados de restrições
cinemáticas é parte central da análise. \\ \addlinespace[4pt]
\autocite{nanograv2024} & Busca correlações escalares transversais; não
encontra evidência significativa a favor de sua inclusão. & Um sinal
correlacionado não basta para identificar uma polarização adicional. \\ \addlinespace[4pt]
\autocite{cordes2025} & Examina funções de redução de sobreposição,
distâncias finitas e corrige um resultado analítico no regime relevante
para gravidade massiva. & É referência de validação; incluir distâncias
finitas não constitui novidade por si só. \\ \addlinespace[4pt]
\autocite{choi2025} & Considera dispersão e polarizações adicionais em
correlações de NANOGrav e CPTA; reporta melhora de ajuste. & Melhor
ajuste não equivale a detecção: é preciso considerar complexidade,
covariâncias e calibração. \\ \addlinespace[4pt]
\autocite{wu2025} & Estuda interferência entre fontes e variabilidade
das correlações no setor tensorial massivo. & A variabilidade entre
realizações do fundo precisa entrar nos testes de robustez. \\ \addlinespace[4pt]
\autocite{zhao2026} & Analisa MeerKAT 4,5 anos com correlações angulares
comprimidas, explora modelos de ruído e faz previsões para SKA. & É o
ponto de partida mais recente encontrado para o recorte principal. \\ \addlinespace[4pt]
\autocite{lvk2026} & Reporta limite de massa e discute efeitos de priori
e de modelagem nos testes de dispersão. & Oferece contexto observacional
e exemplos da necessidade de validar inferências de massa. \\
\end{longtable}
\endgroup
}

O detalhe particularmente útil está na seção III do trabalho de Zhao e
Wang: como os produtos observacionais são comprimidos em frequência, os
autores adotam \(f_{\mathrm{ref}}=1/T\) para relacionar massa e
velocidade. A análise também utiliza uma verossimilhança diagonal por
não dispor da covariância completa dos produtos publicados e limita a
priori pela condição de velocidade real nessa frequência. São hipóteses
explicitadas pelos autores que podem ser investigadas por simulação \autocite{zhao2026}.

A lacuna candidata é \textbf{quantificar, com dados simulados e
calibração estatística, como a compressão em frequência e essas escolhas
afetam limites de massa quando a resposta é dispersiva, incluindo
posteriormente uma contribuição de helicidade zero fisicamente
vinculada}. A busca realizada não estabelece prioridade absoluta sobre
esse recorte. Nos dois primeiros meses deverão ser examinados códigos,
suplementos e trabalhos que citam essas referências.

Também existem apresentações de maio de 2026 de Chris Choi sobre
NANOGrav\footnote{Disponível em:
  \url{https://indico.global/event/16413/contributions/153867/}.
  Consulta em 10 set. 2026.} e de Marcus Bosca e colaboradores sobre
EPTA\footnote{Disponível em:
  \url{https://indico.global/event/16413/contributions/153709/}.
  Consulta em 10 set. 2026.}. Elas indicam pesquisa em andamento e
possível concorrência temática; apresentações de congresso têm status
distinto de artigos revisados por pares.
